% ============================================================
% Introduction section — LNCS style
% Place this between \maketitle / abstract and \section{Related Work}.
% References resolved from references.bib
% ============================================================


\section{Introduction}\label{sec:intro}

Rule-based systems are widely deployed in high-stakes domains such as air traffic control, financial surveillance, and clinical decision support, yet existing benchmarks for production rule engines such as Drools target static, in-memory workloads. No benchmark simultaneously provides: (i)~a \emph{real-world} streaming dataset, (ii)~a domain-derived rule set, (iii)~native temporal CEP operators, and (iv)~a structural characterisation of the rule base. Dayarathna and Perera~\cite{dayarathna2018event} note the absence of agreed-upon streaming benchmarks --- a gap that remains open today.

This paper presents a multi-domain CEP benchmark suite covering financial trading (Binance), collaborative editing (Wikimedia), and aviation surveillance (OpenSky), evaluated on Apache Drools in STREAM mode. Our contributions are:

\begin{enumerate}
  \item \textbf{A benchmark suite spanning three real-world domains}, each comprising a rule set derived from domain standards or operational practice, data collectors for obtaining real-world streaming data from public APIs, and a reproducible replay harness.
  \item \textbf{An empirical cross-suite characterisation} covering static rule-base properties and dynamic runtime metrics, including controlled coverage-variation experiments, demonstrating how domain characteristics --- event arrival regularity, join complexity, and forward-chaining depth --- produce measurable differences in engine behaviour.
\end{enumerate}

Section~\ref{sec:related} surveys related work; Section~3 describes benchmark design; Section~4 presents the cross-suite analysis including coverage-variant experiments; Section~5 concludes.


